%-*-tex-*-
% Copyright Michael J. Ferguson, INRS-Telecommunications
% All rights reserved. 

% Fixed page insert forms to work with pageinsert. ... 22 Sept 87 MJF

% ========= Syntactic Forms ============
% Each of these has a token list that is inserted before each
% particular insert ... and in a group. 
% in \begin... \end.. pairs
\newtoks\everytopinsert
\newtoks\everymidinsert
\newtoks\everypageinsert
\def\begintopinsert{\begingroup\the\everytopinsert\topinsert}
\def\beginmidinsert{\begingroup\the\everymidinsert\midinsert}
\def\beginpageinsert{\begingroup\def\g@box{\vbox to
                         \vinnerpagesize}\def\g@ebox{\relax}\the
                         \everypageinsert\pageinsert}

\def\e@ndinsert{\endinsert\endgroup}

\let\endtopinsert = \e@ndinsert
\let\endmidinsert = \e@ndinsert
\let\endpageinsert = \e@ndinsert

% ============= Figure, Table, etc inserts ==========

% Assumptions 
%   The form is described by three letters, eg "fig" or "tbl"
%   If the list file exists, the page info is written. 
%   \writelistfile must be defined 
%   \autonumber, \autoreference, \ref, \quietref, \veryraggedright
%   If autonumbering is on, then 
%      \silenttrue, \autofignum ..
%   An "fig" form and "tbl" form are defined. These need 
%        \figurelabel, and \tablelabel respectively. These are
%        usually defined in the \englishversion and \versionfrancaise
%   The insert forms use \beginmidinsert ... \endmidinsert. 


%  Defaults  --- to be set. 
%    \captiontextfraction  ... inrsdef.tex 
%    \captionskip  ... inrsdef.tex 
%    \topcaptiontrue or \topcaptionfalse  ... inrsdef.tex 
%    \figurelabel  or \figlabel  ... english.tex or francais.tex 
%    \tablelabel   or \tbllabel  ... english.tex or francais.tex 
%    \captionnumfont ... inrsfont.tex
%    \captiontitlefont ... inrsfont.tex
%    \captionbodyfont ... inrsfont.tex
%    \captionbodyon <off> ... inrsdef.tex


% This is a general form for an  insert with a caption, under it 
% if \topcaptionfalse and over it if \topcaptiontrue
% The particular form is identified with a (usually) three letter sequence
% such as "fig" or "tbl". There can be an associated list file for 
% automatic generation of the page list for the items. 
% The basic form is a non-floating type, enclosed in $$ ... $$ to give
% correct spacing. This can be put in an insert to make it float. 
% Full page forms require \vfill above and/or below for top, bottom, or center.

% The caption box measures its size and centers the text if it is short 
% or puts it in a \vbox which is \captiontextfraction of the \hsize 
% if it is not. 
% The \captionskip is the spacing between the caption and the actual 
% item in the text. This may be a table or figure or ....

% \let\captionnumfont = \bf
% \let\captiontextfont = \relax % defaults to "outside font" 
% \def\captiontextfraction{0.8} 

\def\captionbox#1#2{\setbox0=\hbox{{\captionnumfont #1}\quad 
                       \captiontitlefont #2}
                \ifdim 1\wd0 < \captiontextfraction\hsize
                    \centerline{\unhbox0}
                \else
                    \centerline{\vtop{\hsize=\captiontextfraction\hsize
                                      \veryraggedright
                                      \setbox0=\hbox{{\captionnumfont #1}\quad}
                                      \hangindent=1\wd0 \hangafter=1
                                      \noindent\unhbox0 
                                        \captiontitlefont #2}}%
                \fi}

% A caption is assumed to have both a <title> that will appear in the 
% list, if any and a <caption body> that will appear only when the caption is
% printed. The default is to have \captionbody ... "active" and to turn it
% off when making lists. 
% This is done with two commands \captionbodyon and \captionbodyoff 
% associated with a <caption body> is a <caption body font>

\def\captionbodyon{\def\captionbody##1{{\captionbodyfont ##1}}}
\def\captionbodyoff{\def\captionbody##1{}}


%  \rect <hor dimen>  by <vert dimen>}  A centered rectangular box to act 
% as a place holder if there is a figure to be pasted in later. 
\def\r@ect{\centerline{\vbox{\hrule
                          \hbox{\vrule height \the\dimen2
                                \hskip \dimen0
                                \vrule height \the\dimen2}
                          \hrule
                         }}}
\def\rect #1by{\dimen0 #1\afterassignment\r@ect \dimen2}

% ======= This is the general insert form =======
% #1 is the name of the insert form, #2 is the tag or number, #3 is the title
% #4 is the actual insert form value ... ie figure or table or ...

% This is the glue specification between the caption and the body of the 
% insert

\newskip\captionskip  %\captionskip=2ex

% This is a switch that controls whether the caption is above or below the
% the body of the insert

\newif\iftopcaption  %\topcaptionfalse % defaulted 

% This is the general form of the insert --- it requires both \autonumber
% and \autoreference to work. By selectively turning on/off these
% switches, multipart figures are possible. 

% \ifsilentfigure ... turns off the figure inserts so that tagging can 
% occur with figures gracefully occuring at the end. This means that 
% figure list file will refer to pages where the figures are sort of
% called out 

\def\multiinsert{\def\g@box{\vbox }\def\g@ebox{\relax}}

\newif\ifsilentfigure  % default is \silenfigurefalse 

\def\g@box{$$\vbox } \def\g@ebox{$$}

\def\geninsertform #1#2#3#4{
                    \ifautonumber 
                  {\silenttrue\expandafter\csname auto#1num\endcsname{#2}}\relax
                        \edef\i@nsnum{\csname:@#2\endcsname}\else
                        \edef\i@nsnum{#2}\fi 
                        \writelistfile{#1}{\csname #1tocout\endcsname}{\i@nsnum}{#3}
                  \ifsilentfigure \relax \else
                   \g@box{\normalbaselines
                   \iftopcaption
                         \captionbox{\csname #1label\endcsname\i@nsnum}{#3}
                         \vskip \captionskip
                         \centerline{\vbox{#4}}
                     \else
                         \centerline{\vbox{#4}}
                         \vskip \captionskip
                         \captionbox{\csname #1label\endcsname\i@nsnum}{#3}
                     \fi
                    \vss}\g@ebox\fi}


%  ==========  Special Cases ============
%  These define simple forms for standard midinserts. It uses 
%  \beginmidinsert ... \endmidinsert for consistency

% For Figures 
\def\figureform #1#2#3{\geninsertform {fig}{#1}{#2}{#3}}
\def\figlabel{\figurelabel }
\def\figureinsert #1#2#3{\beginmidinsert 
                         \figureform{#1}{#2}{#3} 
                         \endmidinsert }

%\def\figtocout{\string\figtoc} % actually set in inrsdef.tex 


% For Tables 
\def\tableform #1#2#3{\geninsertform {tbl}{#1}{#2}{#3}}
\def\tbllabel{\tablelabel }
\def\tableinsert #1#2#3{\beginmidinsert 
                         \tableform{#1}{#2}{#3} 
                         \endmidinsert }
%\def\tbltocout{\string\tbltoc} % actually set in inrsdef.tex 


% Comment ---- \endinsert in Plain has a pre and post \bigskip 
% This should probably be removed and the spacing always controlled by 
% the $$ ... $$. Instead, if the \everymidinsert form were used it 
% would be possible to remove it with \everymidinsert = {\let\bigskip=\relax}